% !TEX root = ../main.tex
%% Sujet type X-ENS — Inégalités de Hardy et de Carleman.
%% Résultats vérifiés par verifications/eval_cc.py.

\sujetEval{C4}{Sujet type X-ENS · Mathématiques}{Sujet type X-ENS}
  {Durée : \textbf{4 heures}\\Calculatrice interdite\\Inégalités de moyennes}{4}
  {Partie A/L'inégalité de Hardy/Q1 à Q4,Partie B/La constante 4 est optimale/Q5 et Q6,%
   Partie C/L'inégalité de Carleman/Q7 à Q10,Partie D/La constante $\mathrm{e}$ est optimale/Q11 et Q12}
\enParties
\consigne{Sujet inédit, conçu dans l'esprit des épreuves de l'X et des ENS :
peu de calculs, mais des raisonnements à construire. On pourra utiliser
l'inégalité arithmético-géométrique (devoir maison \numero1) et l'inégalité
$\bigl(1+\frac1k\bigr)^{k}<\mathrm{e}$, valable pour $k\geqslant1$.}

Dans tout le problème, $(a_n)_{n\geqslant1}$ est une suite de réels
positifs, et l'on pose, pour $n\geqslant1$,
\[A_n=\frac{a_1+a_2+\dots+a_n}{n}\qquad\text{et}\qquad G_n=\bigl(a_1a_2\cdots a_n\bigr)^{1/n},\]
ainsi que $A_0=0$. On compare les séries $\sum A_n^{2}$ et $\sum a_n^{2}$,
puis $\sum G_n$ et $\sum a_n$.

\begin{exoEval}[L'inégalité de Hardy]{6 points}{}
\begin{questions}
  \item Montrer que $a_n=nA_n-(n-1)A_{n-1}$ pour tout $n\geqslant1$.
  \item En déduire que
    $A_n^{2}-2a_nA_n\leqslant(n-1)A_{n-1}^{2}-nA_n^{2}$ pour $n\geqslant1$.
  \item Soit $N\geqslant1$. Montrer que
    $\sum_{n=1}^{N}A_n^{2}\leqslant2\sum_{n=1}^{N}a_nA_n$, puis que
    \[\sum_{n=1}^{N}A_n^{2}\leqslant4\sum_{n=1}^{N}a_n^{2}.\]
  \item En déduire que si $\sum a_n^{2}$ converge, alors $\sum A_n^{2}$
    converge et $\sum_{n=1}^{+\infty}A_n^{2}\leqslant4\sum_{n=1}^{+\infty}a_n^{2}$
    (\emph{inégalité de Hardy}).
\end{questions}
\end{exoEval}

\begin{exoEval}[La constante 4 est optimale]{4 points}{}
On fixe $N\geqslant1$ et l'on choisit $a_n=\frac{1}{\sqrt n}$ pour
$n\leqslant N$, $a_n=0$ pour $n>N$.
\begin{questions}[start=5]
  \item Montrer que pour $n\geqslant1$ :
    $2\sqrt{n+1}-2\leqslant\sum_{k=1}^{n}\dfrac{1}{\sqrt k}\leqslant2\sqrt n-1$.
  \item En déduire qu'il existe une constante $K$, indépendante de $N$,
    telle que $\sum_{n=1}^{N}A_n^{2}\geqslant4H_N-K$, où
    $H_N=\sum_{n=1}^{N}\frac1n$. Montrer que $4$ ne peut être remplacé, dans
    l'inégalité de Hardy, par une constante strictement plus petite.
\end{questions}
\end{exoEval}

\begin{exoEval}[L'inégalité de Carleman]{6 points}{}
Pour $k\geqslant1$, on pose $c_k=\dfrac{(k+1)^{k}}{k^{k-1}}$.
\begin{questions}[start=7]
  \item Montrer que $c_1c_2\cdots c_n=(n+1)^{n}$ pour tout $n\geqslant1$.
  \item En appliquant l'inégalité arithmético-géométrique aux réels
    $a_1c_1,\dots,a_nc_n$, montrer que
    $G_n\leqslant\dfrac{1}{n(n+1)}\sum_{k=1}^{n}a_kc_k$.
  \item Montrer que $\sum_{n=k}^{N}\dfrac{1}{n(n+1)}\leqslant\dfrac1k$, puis que
    \[\sum_{n=1}^{N}G_n\leqslant\sum_{k=1}^{N}a_k\Bigl(1+\frac1k\Bigr)^{k}.\]
  \item En déduire que si $\sum a_n$ converge, alors $\sum G_n$ converge et
    $\sum_{n=1}^{+\infty}G_n\leqslant\mathrm{e}\sum_{n=1}^{+\infty}a_n$
    (\emph{inégalité de Carleman}).
\end{questions}
\end{exoEval}

\begin{exoEval}[La constante $\mathrm{e}$ est optimale]{4 points}{}
\begin{questions}[start=11]
  \item Montrer que pour $n\geqslant1$ :
    $n\ln n-n+1\leqslant\ln n!\leqslant(n+1)\ln(n+1)-n$. En déduire que
    $(n!)^{1/n}\sim\dfrac{n}{\mathrm{e}}$.
  \item On fixe $N\geqslant1$ et l'on choisit $a_n=\frac1n$ pour
    $n\leqslant N$, $a_n=0$ pour $n>N$. Soit $\varepsilon\in\intOO01$.
    Montrer qu'il existe $n_0$, indépendant de $N$, tel que
    $\sum_{n=1}^{N}G_n\geqslant(1-\varepsilon)\,\mathrm{e}\,(H_N-H_{n_0})$
    pour $N\geqslant n_0$. Conclure.
\end{questions}
\end{exoEval}
